% Extension of AG.39; intended for incorporation into that statement/proof,
% not a second independently numbered configuration claim.

\paragraph{The congruence lemma needed for every excluded prime.}
Let \(\Gamma\subset\operatorname{GL}_N(\mathbb Z[1/M])\) be
finitely generated. Suppose that the identity component \(G\) of its
Zariski closure over \(\mathbb Q\) is a nontrivial semisimple group.
For every prime \(p\) and every integer \(a\geq1\), there are
infinitely many rational primes \(\ell\ne p\), with \(\ell\nmid M\),
such that
\[
       p^a\mid
       \#\operatorname{im}\bigl(\Gamma\longrightarrow
                    \operatorname{GL}_N(\mathbb F_\ell)\bigr).
                                                        \tag{AG.39.SA}
\]

\begin{proof}
The subgroup \(\Gamma_0=\Gamma\cap G(\mathbb Q)\) has finite index,
is finitely generated, and is Zariski dense in \(G\). Let
\[
     1\longrightarrow Z\longrightarrow \widetilde G
          \xrightarrow{\pi}G\longrightarrow1
\]
be its simply connected central covering, and put \(c=\#Z\).
The connecting map
\[
               \delta:G(\mathbb Q)\longrightarrow H^1(\mathbb Q,Z)
\]
is a group homomorphism. Indeed, for lifts of two rational elements,
the two Galois cocycles multiply because their values lie in the
central subgroup \(Z\). The abelian group \(H^1(\mathbb Q,Z)\)
is killed by \(c\). Therefore the image of the finitely generated
group \(\Gamma_0\) under \(\delta\) is finite. Its kernel
\(\Gamma_1\) has finite index in \(\Gamma_0\), and every element
of \(\Gamma_1\) has a rational lift through \(\pi\).

Choose a finite generating set of \(\Gamma_1\), choose rational
lifts of its generators, and let \(\widetilde\Gamma\) be the
group generated by those lifts. Its image is exactly \(\Gamma_1\).
Its Zariski closure is all of \(\widetilde G\): the image of that
closure under the finite morphism \(\pi\) is closed and contains
the dense subgroup \(\Gamma_1\); equality of dimensions, and
connectedness of \(\widetilde G\), finish the assertion.

Choose a faithful rational representation of \(\widetilde G\).
After enlarging a finite set of primes \(S\), all chosen lifts
are integral over \(\mathbb Z[S^{-1}]\), and \(\pi\), the
original faithful representation of \(G\), and their central
kernel spread out over this ring. We may require that these
models are smooth semisimple group schemes, that the morphism
is a central isogeny of degree \(c\), and that the embedding
of \(G\) into the original matrix group is a closed immersion.
Thus, outside \(S\), reducing \(\pi(\widetilde\gamma)\) gives
the same matrix as reducing the original element of \(\Gamma_1\).
The kernel of this comparison on finite-field points has at
most \(c\) elements.

Nori's strong-approximation theorem, in precisely its
connected and simply connected form
\cite[Theorem 5.4, printed p. 274]{AC_Nori1987}, says that the
congruence closure of \(\widetilde\Gamma\) is open in
\(\widetilde G(\prod_{\ell\notin S}\mathbb Z_\ell)\).
Consequently, after excluding finitely many more primes,
\[
       \widetilde\Gamma\bmod\ell
                      =\widetilde G(\mathbb F_\ell).
\]
Here smoothness gives surjectivity of reduction from
\(\mathbb Z_\ell\)-points, and an open subgroup of a product
contains the full product outside finitely many factors.

Choose a finite Galois number field \(K\) which splits
\(\widetilde G\) and contains the roots of unity of order
\(p^{a+v_p(c)}\). Exclude also the finitely many rational
primes needed to spread the chosen splitting isomorphism
over the ring of integers of \(K\). Infinitely many rational primes split
completely in \(K\). One may use the usual splitting-prime
theorem, or the following elementary sufficient form. If
\(K=\mathbb Q(\alpha)\) with \(\alpha\) integral and minimal
polynomial \(f\), infinitely many primes divide nonzero values
\(f(n)\). Indeed, if all such prime divisors belonged to a
finite set with product \(R\), choose \(b\) with
\(d=f(b)\ne0\). The integer polynomial
\(f(b+dRt)/d\) is congruent to one modulo \(R\) and has
nonconstant absolute value, a contradiction for sufficiently
large integers \(t\). Away from the discriminant and the
index of \(\mathbb Z[\alpha]\), a root of \(f\) modulo a
prime gives a residue-degree-one prime of \(K\). Since
\(K/\mathbb Q\) is Galois, that rational prime splits completely.

Discard the finite exceptional set, and choose any such
\(\ell\ne p\). Then \(\widetilde G_{\mathbb F_\ell}\)
is a nontrivial split semisimple group and
\(p^{a+v_p(c)}\mid\ell-1\). Its split maximal torus has
positive rank, so
\[
       v_p\bigl(\#\widetilde G(\mathbb F_\ell)\bigr)
                          \geq a+v_p(c).
\]
Its image in the original \(\operatorname{GL}_N(\mathbb F_\ell)\)
loses at most \(v_p(c)\) from this valuation, and that image is
the reduction of \(\Gamma_1\). It is a subgroup of the
reduction of \(\Gamma\). This proves (AG.39.SA).
\end{proof}

\paragraph{Replacement of the pro-odd restriction in AG.39.}
Fix an arbitrary prime \(p\), and write \((-)^{(p')}\) for
completion with respect to all finite groups of order prime
to \(p\). For every smooth proper complex curve \(C\) of
genus \(g\geq2\), every \(P\in C\), and
\(U=(C\times C)-\Delta_C\), the map
\[
       \pi_1(C-\{P\})^{(p')}
                 \longrightarrow \pi_1(U)^{(p')}
                                                        \tag{AG.39.ALL}
\]
is not injective. The corresponding geometric prime-to-\(p\)
map is likewise noninjective for every smooth proper curve of
genus at least two over an algebraically closed field of
characteristic \(p\).

\begin{proof}
In the explicit family construction of AG.39 choose a prime
\(r\ne p\): one may take \(r=2\) for odd \(p\), and
\(r=3\) for \(p=2\). Replace its cyclic triple cover by a
cyclic degree-\(r\) etale cover \(a:C'\to C\), and take
\(b:B\to C\) to be the full mod-\(r\) homology cover,
of degree \(r^{2g}\). This kills the cross term of the
diagonal class modulo \(r\), and both remaining Kunneth
terms vanish modulo \(r\). The same integral-cohomology and
Picard-divisibility argument therefore gives
\[
       L^{\otimes r}\simeq
       \mathcal O_{C'\times B}\bigl((a\times b)^{-1}\Delta_C\bigr).
\]
The resulting cyclic cover is a smooth proper family
\(S\to B\). Each fiber is a degree-\(r\) cover of \(C'\)
branched at the \(r\) distinct points \(a^{-1}(b(z))\), and
\[
   g(C')=1+r(g-1),\qquad g(B)=1+r^{2g}(g-1),\qquad
   g(S_z)=1+r^2(g-1)+\frac{r(r-1)}2.
\]
It is nonisotrivial by the same finite-Hom argument: a fixed
source curve has only finitely many degree-\(r\) maps to the
fixed hyperbolic curve \(C'\), while their required branch
sets vary through the fibers of the surjection \(a\).

Every cover in the punctured total-space tower has an
\(r\)-group Galois group, and its composite has an
\(r\)-group Galois closure by the wreath-product argument.
The punctured fiber has the same property. An open subgroup
containing a normal subgroup of finite prime-to-\(p\) index
has its full pro-prime-to-\(p\) topology induced from the
ambient group: taking normal cores produces an extension
of a prime-to-\(p\) group by a subgroup of a product of such
groups. This uses no solvability assumption.

Consequently, if (AG.39.ALL) were injective, the completed
punctured-fiber group would be a closed normal subgroup of
the completed total group of \(S^\circ\to B\). For every
\(\ell\ne p\), its compactified homology
\(H_1(S_z,\mathbb F_\ell)\) is a finite quotient with
monodromy-invariant kernel. Conjugation is continuous, the
fiber acts trivially on this abelian quotient, and the
discrete total group surjects onto \(\pi_1(B)\). Therefore
every reduction modulo \(\ell\ne p\) of the integral
monodromy group
\[
       \Gamma\subset\operatorname{Sp}_{2g(S_z)}(\mathbb Z)
\]
would have order prime to \(p\). These are reductions of
the actual integral homology local system, rather than
abstract finite quotients of an unrelated group.

Deligne's semisimplicity theorem used in AG.39 makes the
connected Zariski closure of \(\Gamma\) semisimple. It is
nontrivial: a trivial identity component would make
\(\Gamma\) finite, so after a finite etale base change the
period map would be a holomorphic map from a compact curve
to the bounded realization of Siegel space. It would be
constant; Torelli would make the family isotrivial. The
congruence lemma, with \(a=1\), now contradicts the asserted
prime-to-\(p\) orders of all the reductions. This proves
(AG.39.ALL) over \(\mathbb C\).

For characteristic \(p\), lift the pointed smooth curve
as in AG.39 and apply the prime-to-\(p\) specialization
isomorphisms for the two smooth proper pairs
\((\mathcal C,\mathcal P)\) and
\((\mathcal C\times\mathcal C,\Delta_{\mathcal C})\).
SGA~1 XIII, Corollary~2.9 and Section~2.10 apply because
all the finite groups in question have order invertible
on the base. Functoriality identifies the two fiber maps.
Thus injectivity of the special-fiber map would imply
injectivity of its complex counterpart, which has just
been contradicted. The Hodge and period-map argument has
been used only in characteristic zero.
\end{proof}

The extension eliminates the full configuration exactness
shortcut for every excluded prime. It does not produce a
collision, and it leaves A1 and universal relative-witness
existence unresolved. The characteristic-two proof in AG.39
via Feit--Thompson and bounded solvable length remains an
independent alternative to the strong-approximation step.
